\begin{table}[!htbp]
\centering
\sbox{\tabletempbox}{%
\footnotesize
\begin{tabular}[t]{rrrr>{\raggedright\arraybackslash}p{15em}}
\toprule
Percentile & Actual $t(\hat{\alpha})$ & Simulated Mean & Prob.\ Luck & Interpretation\\
\midrule
1 & -2.477 & -2.882 & 79.2\% & Consistent with zero skill\\
5 & -1.871 & -1.889 & 48.5\% & Consistent with zero skill\\
10 & -1.352 & -1.434 & 57.1\% & Consistent with zero skill\\
50 & 0.019 & 0.018 & 50.5\% & Consistent with zero skill\\
90 & 1.559 & 1.475 & 62.6\% & Consistent with zero skill\\
95 & 2.047 & 1.940 & 64.0\% & Consistent with zero skill\\
99 & 2.742 & 2.994 & 34.5\% & Consistent with zero skill\\
\bottomrule
\end{tabular}%
}
\setlength{\tabletempwidth}{\wd\tabletempbox}
\ifdim\tabletempwidth>\linewidth\setlength{\tabletempwidth}{\linewidth}\fi
\begin{minipage}{\tabletempwidth}
\captionsetup{width=\linewidth}
\caption{\label{tab:bootstrap_tails_FF}Fama--French (2010) Bootstrap: Actual vs.\ Simulated $t(\hat{\alpha})$ Percentiles -- FF Subperiod}
\ifdim\wd\tabletempbox>\linewidth
  \resizebox{\linewidth}{!}{\usebox{\tabletempbox}}
\else
  \usebox{\tabletempbox}
\fi
\par\medskip
\begin{singlespace}\footnotesize\noindent
Bootstrap procedure follows Fama and French (2010, \textit{Journal of Finance}). Sample: actively managed funds, Jan 1995--Sep 2006 (Fama-French 2010 subperiod), minimum 24 monthly observations ($N = 804$ funds). For each fund, estimated monthly alpha is subtracted from the excess return series to construct a zero-alpha null return. In each of $B = 10{,}000$ bootstrap iterations, calendar months are resampled with replacement and kept in the order drawn, the same draw applying to all funds, preserving cross-sectional factor return dependence. The Carhart (1997) four-factor model is re-estimated on each resampled series. \textit{Actual} $t(\hat{\alpha})$: percentile of the empirical $t$-statistic distribution across active funds. \textit{Simulated Mean}: average of that percentile across all iterations. \textit{Prob.\ Luck}: fraction of iterations in which the simulated percentile falls below the actual value; values below 5\% (above 95\%) indicate that the actual percentile is worse (better) than zero-alpha luck would produce; 0.0\% means no iteration and $<$0.1\% means 1 to 4 of 10,000 iterations. Newey-West standard errors with a 6-month lag are used throughout. Compare with Fama and French (2010), Table III.
\end{singlespace}
\end{minipage}
\end{table}

